\documentclass[12pt,a4paper]{article}
\usepackage[UTF8,heading=true]{ctex}
\usepackage{geometry}
\usepackage{amsmath,amssymb}
\usepackage{booktabs,array,tabularx}
\usepackage{enumitem,hyperref,fancyhdr,xcolor,setspace}
\usepackage{ragged2e}

\newcolumntype{L}[1]{>{\RaggedRight\arraybackslash}p{#1}}
\ctexset{
    section = {
        name = {,、},
        number = \chinese{section},
        format = {\Large\bfseries\raggedright}
    }
}

\geometry{left=2.5cm,right=2.5cm,top=2.5cm,bottom=2.5cm}
\setlength{\headheight}{15pt}
\setstretch{1.5}
\setlength{\emergencystretch}{3em}
\hfuzz=20pt
\sloppy
\hypersetup{hidelinks}

\pagestyle{fancy}
\fancyhf{}
\fancyhead[C]{权利要求书}
\fancyfoot[C]{\thepage}
\renewcommand{\headrulewidth}{0pt}

\begin{document}

\begin{center}
    {\Huge\heiti 中国移动专利申请} \\[0.5em]
    {\Huge\heiti 权利要求书} \\[1em]
\end{center}

\noindent
发明名称：长距离量子中继器中的分布式纠缠纯化架构

\vspace{0.5em}

\begin{enumerate}[label=\arabic*.,leftmargin=2em,itemsep=0.6em]
    \item 一种用于量子中继网络的分布式纠缠纯化与纠缠交换的方法，其特征在于，包括：
    在相邻的量子中继节点之间建立基础纠缠对，并为所述基础纠缠对分配纠缠对标识；
    在每一链路段中将最先满足建链确认条件的纠缠对登记为锚定纠缠对，并将随后生成的纠缠对登记为候选牺牲纠缠对；
    对所述锚定纠缠对和所述候选牺牲纠缠对获取或估计至少包括保真度与年龄的状态信息，并基于所述状态信息计算综合评分以生成决策；
    当所述综合评分满足纯化触发条件且所述锚定纠缠对的年龄未超过丢弃阈值时，在与相邻节点协同的纯化窗口内，对所述锚定纠缠对与所述候选牺牲纠缠对执行双边门序列与校验测量，并依据校验结果更新所述锚定纠缠对的状态或丢弃本轮涉及的局部资源；
    当所述锚定纠缠对满足交换门限且位于某中间节点的左右两侧锚定纠缠对均满足交换条件时，在交换窗口内于所述中间节点执行贝尔态测量以完成纠缠交换，并将交换结果通知至相关端点节点以更新长距离纠缠对的状态；
    在经典信道中传输压缩综合征消息，以使节点间仅依据决策充分信息完成所述纯化触发、交换触发、资源释放与状态更新。

    \item 根据权利要求1所述的方法，其特征在于，所述纠缠对标识包括链路段标识、层级标识、生成轮次和时间戳中的至少两项，其中所述层级标识用于表示基础纠缠对以及由纠缠交换递归形成的更高层级纠缠对。

    \item 根据权利要求1所述的方法，其特征在于，所述综合评分至少由锚定纠缠对的保真度估计、锚定纠缠对的年龄、候选牺牲纠缠对的保真度估计以及链路重建代价的估计共同确定，并用于在纯化、交换与丢弃之间进行选择。

    \item 根据权利要求1所述的方法，其特征在于，在一次建链尝试成功后，将由通信量子比特持有的纠缠态通过局部门操作写入存储量子比特以释放通信量子比特继续执行下一次建链尝试，从而实现通信量子比特与存储量子比特的资源复用。

    \item 根据权利要求1所述的方法，其特征在于，所述纯化窗口与所述交换窗口被抽象为可重叠的操作窗口，且在窗口内对参与本轮操作的量子比特执行锁定以避免冲突操作，在窗口外释放所述量子比特用于建链或后续操作。

    \item 根据权利要求5所述的方法，其特征在于，所述压缩综合征消息中携带租约字段，所述租约字段至少包括会话标识与有效时间，租约到期时自动释放对应的量子资源锁定以降低死锁风险。

    \item 根据权利要求1所述的方法，其特征在于，所述丢弃阈值与量子存储相干时间常数相关联，当锚定纠缠对年龄超过所述丢弃阈值时触发局部丢弃与重建，以降低退相干导致的无效等待。

    \item 根据权利要求1所述的方法，其特征在于，所述压缩综合征消息至少包括锚定纠缠对标识、最近一轮纯化或交换的成功失败标志、必要的校验测量结果、时间戳或年龄信息以及更新后的保真度估计，其中用于纠缠交换的贝尔态测量结果采用二比特表示以用于帕里框架更新。

    \item 根据权利要求1所述的方法，其特征在于，所述压缩综合征消息进一步包括单调递增计数器以用于重放抑制，且节点仅基于所述计数器与所述纠缠对标识对重复消息进行丢弃或合并处理。

    \item 根据权利要求1所述的方法，其特征在于，所述方法将失败恢复限定在当前链路段或当前会话范围内，在纯化失败或交换失败时仅回退受影响的局部资源而不触发全局回退。

    \item 一种量子中继节点装置，其特征在于，包括：
    通信量子比特，用于与相邻节点建立基础纠缠对；
    存储量子比特，用于缓存锚定纠缠对和候选牺牲纠缠对；
    辅助量子比特，用于在纯化窗口内执行双边门序列与校验测量，或在交换窗口内辅助控制；
    光子接口或光子芯片子模块，用于纠缠对制备、贝尔态测量和经典结果上报；
    本地控制器，用于生成纠缠对标识，执行综合评分计算，维护状态机并驱动纯化与纠缠交换的窗口化调度；
    其中，所述本地控制器被配置为执行如权利要求1至10任一项所述的方法。

    \item 根据权利要求11所述的装置，其特征在于，所述本地控制器包括候选牺牲纠缠对管理单元，用于对新生成纠缠对按综合评分进行登记、替换与回收，从而以滚动方式为锚定纠缠对提供牺牲纠缠对。

    \item 根据权利要求11所述的装置，其特征在于，所述本地控制器包括窗口锁定单元，用于在纯化窗口或交换窗口内对参与操作的量子比特建立锁定并在窗口结束后释放。

    \item 根据权利要求11所述的装置，其特征在于，所述本地控制器包括租约管理单元，用于在跨节点协同操作前生成租约并写入压缩综合征消息，且在租约到期或会话完成时撤销锁定。

    \item 根据权利要求11所述的装置，其特征在于，所述本地控制器包括年龄监测单元，用于基于时间戳计算锚定纠缠对年龄并在超过丢弃阈值时触发局部丢弃与重建。

    \item 根据权利要求11所述的装置，其特征在于，所述本地控制器包括综合征压缩单元，用于将纯化或交换过程中产生的原始测量记录压缩为包含成功失败标志、必要校验比特与帕里框架更新信息的压缩综合征消息。

    \item 一种量子中继系统，其特征在于，包括多个量子中继节点装置以及与所述多个量子中继节点装置通信连接的经典信道，其中所述量子中继节点装置为权利要求11至17任一项所述的量子中继节点装置，所述量子中继系统被配置为在异步链路条件下基于压缩综合征消息实现分布式纠缠纯化与纠缠交换的流水线调度。

    \item 一种电子设备，其特征在于，包括处理器与存储器，所述存储器中存储有计算机程序，所述计算机程序被所述处理器执行时用于实现权利要求1至10任一项所述的方法中的经典控制与调度逻辑。

    \item 一种计算机可读存储介质，其特征在于，其上存储有计算机程序，所述计算机程序被处理器执行时用于实现权利要求1至10任一项所述的方法中的经典控制与调度逻辑。
\end{enumerate}

\end{document}

